% Figure: HPA Axis Model with ME/CFS Dysfunction
% Mathematical model diagram for Chapter 29 (Neuroendocrine Models)
% Shows cascade, feedback loops, inputs, and ME/CFS dysfunction annotations

\begin{figure}[htbp]
\centering
\resizebox{\textwidth}{!}{%
\begin{tikzpicture}[
    node distance=2.5cm and 3cm,
    >=Stealth,
    every node/.style={font=\sffamily},
    % Organ boxes
    organ/.style={
        draw=orange!70!black, fill=orange!10, very thick,
        rounded corners, text width=4.2cm, align=center,
        minimum height=1.2cm, font=\sffamily\bfseries
    },
    % Hormone labels on arrows
    hormone/.style={
        font=\sffamily\small, fill=white, inner sep=2pt
    },
    % Negative feedback arrows
    inhibit/.style={
        -Bar, thick, red!70!black, dashed
    },
    % Input arrows
    input-arrow/.style={
        -Stealth, thick, blue!70!black
    },
    % Input source box
    input-box/.style={
        draw=blue!60!black, fill=blue!5, rounded corners,
        font=\sffamily\small, text=blue!70!black,
        text width=2.8cm, align=center, minimum height=0.8cm,
        inner sep=3pt
    },
    % Input equation label
    input-eq/.style={
        font=\sffamily\scriptsize, text=blue!50!black
    },
    % ME/CFS dysfunction annotation
    mecfs/.style={
        draw=red!70!black, fill=red!8, rounded corners,
        font=\sffamily\scriptsize, text=red!70!black,
        inner sep=4pt, align=left
    },
    % Connection to immune system
    immune-box/.style={
        draw=purple!70!black, fill=purple!8, very thick,
        rounded corners, text width=3cm, align=center,
        minimum height=1cm, font=\sffamily\bfseries
    },
    % Cascade arrow
    cascade/.style={
        -Stealth, very thick, orange!70!black
    },
]

% ============================================================
% TITLE
% ============================================================
\node[font=\sffamily\large\bfseries, orange!70!black]
    (title) at (1, 1.5) {HPA Axis Model with ME/CFS Dysfunction};

% ============================================================
% MAIN CASCADE (center column)
% ============================================================

\node[organ] (hypo) at (1, 0) {Hypothalamus};
\node[organ, below=2.8cm of hypo] (pit) {Anterior Pituitary};
\node[organ, below=2.8cm of pit] (adr) {Adrenal Cortex};
\node[organ, below=2.8cm of adr, draw=orange!50!black, fill=orange!5]
    (target) {Target Tissues / Effects};

% ============================================================
% CASCADE ARROWS with hormone labels
% ============================================================

\draw[cascade] (hypo.south) -- (pit.north)
    node[hormone, midway, left=4pt] {CRH ($H$)};
\draw[cascade] (pit.south) -- (adr.north)
    node[hormone, midway, left=4pt] {ACTH ($A$)};
\draw[cascade] (adr.south) -- (target.north)
    node[hormone, midway, left=4pt] {Cortisol ($F$)};

% ============================================================
% NEGATIVE FEEDBACK LOOPS (right side)
% Feedback taps off the cortisol arrow (midway adr→target)
% Routes rightward then up, entering organ boxes from the right
% ============================================================

% Feedback tap point on cortisol arrow
\coordinate (fb-tap) at ($ (adr.south)!0.5!(target.north) $);

% --- Long loop: Cortisol → Hypothalamus ---
\draw[inhibit]
    (fb-tap) -- ++(5cm, 0) coordinate (fb-long-corner)
    |- ($ (hypo.east) + (0, 0) $);
\node[hormone, text=red!70!black, anchor=west] at (6.1, -2.5) {$K_F,\; n_F$};

% --- Short loop: Cortisol → Pituitary ---
\draw[inhibit]
    (fb-tap) -- ++(3.5cm, 0) coordinate (fb-short-corner)
    |- ($ (pit.east) + (0, 0) $);
\node[hormone, text=red!70!black, anchor=west] at (4.6, -4.5) {$K_F,\; n_{FA}$};

% Feedback annotation
\node[font=\sffamily\scriptsize, text=red!60!black, align=center]
    at (6.1, 0.8) {Negative\\feedback};

% ============================================================
% INPUTS (left side)
% ============================================================

% Circadian clock
\node[input-box] (circ) at (-3.5, 0.4)
    {Circadian clock};
\node[input-eq, below=1pt of circ]
    {$\sin(2\pi t/24 - \varphi_c)$};
\draw[input-arrow] (circ.east) -- (hypo.west |- circ.east);

% Stress
\node[input-box] (stress) at (-3.5, -0.9)
    {Stress input};
\node[input-eq, below=1pt of stress]
    {$\sigma_{\text{stress}}(t)$};
\draw[input-arrow] (stress.east) -- (hypo.west |- stress.east);

% Immune cytokines
\node[input-box] (cyto) at (-3.5, -2.2)
    {Immune cytokines};
\node[input-eq, below=1pt of cyto]
    {from immune model};
\draw[input-arrow] (cyto.east) -- ++(0.5, 0) |- ($ (hypo.south west) + (0, 0.1) $);

% ============================================================
% ME/CFS DYSFUNCTION ANNOTATIONS
% ============================================================

% At Hypothalamus: enhanced feedback sensitivity (right, near feedback)
\node[mecfs, text width=3.5cm, anchor=west] (mecfs-hypo) at (7, -1.2)
    {$\uparrow\; n_F$ enhanced\\feedback sensitivity};
\draw[red!50!black, thick, dotted, -Stealth]
    (mecfs-hypo.north) -- ++(0, 0.5) -| (fb-long-corner);

% At Circadian: blunted amplitude (left of circadian box)
\node[mecfs, text width=3cm, anchor=east] (mecfs-circ) at (-6.2, 0.4)
    {$\downarrow\; a_c$ blunted\\circadian amplitude};
\draw[red!50!black, thick, dotted, -Stealth]
    (mecfs-circ.east) -- (circ.west);

% At Stress: reduced gain (left of stress box)
\node[mecfs, text width=3cm, anchor=east] (mecfs-stress) at (-6.2, -0.9)
    {$\downarrow\; \sigma_{\text{stress}}$ gain};
\draw[red!50!black, thick, dotted, -Stealth]
    (mecfs-stress.east) -- (stress.west);

% At Cortisol output: low-normal (right, at target tissues level)
\node[mecfs, text width=3.5cm, anchor=west]
    (mecfs-cort) at (7, 0 |- target) {Low-normal cortisol\\output};
\draw[red!50!black, thick, dotted, -Stealth]
    (mecfs-cort.west) -- (target.east);

% ============================================================
% CONNECTION TO IMMUNE SYSTEM (bottom)
% ============================================================

\node[immune-box, below=2cm of target] (brake)
    {Anti-inflammatory\\Brake};
\node[immune-box, right=3cm of brake] (immune)
    {Immune System\\{\normalfont\small(Eq.~\ref{eq:cortisol-immune})}};

\draw[cascade] (target.south) -- (brake.north)
    node[hormone, midway, left=4pt] {$F$};
\draw[inhibit] (brake.east) -- (immune.west)
    node[hormone, midway, above=4pt, text=purple!70!black] {suppresses};

% ME/CFS annotation at immune connection
\node[mecfs, below=0.8cm of brake, text width=7.5cm, align=center]
    (mecfs-immune) {Weakened in ME/CFS
    $\rightarrow$ permits sustained inflammation};
\draw[red!50!black, thick, dotted]
    (mecfs-immune.north) -- (brake.south);

\end{tikzpicture}%
}% end resizebox
\caption{HPA axis feedback model (Equation~\ref{eq:hpa-axis}). The
hypothalamus--pituitary--adrenal cascade produces cortisol under circadian and
stress-driven regulation with negative feedback at hypothalamic and pituitary
levels. In ME/CFS (red annotations), enhanced feedback sensitivity, blunted
circadian amplitude, and reduced stress responsiveness produce low-normal
cortisol that inadequately restrains immune activation.}
\label{fig:hpa-axis-model}
\end{figure}
